\chapter{هندسة اللف المغزلي والحالة المتشابكة الأولى}

\lectureoverview{ينتمي التشابك إلى الأنظمة المؤلفة من أجزاء، لكن لا يمكن تناول مسألة اللفين المغزليين بأمانة ما لم تكتمل أولًا أدوات اللف المغزلي الواحد. لذلك نبدأ باللف المغزلي على امتداد محور اعتباطي، ونشتق احتمال الانتقال بين اتجاهين، ونفسر القياس بوصفه مرشحًا، ونحدد التبادلية بوصفها معيارًا للمرصودات المتوافقة. ثم تكشف حجة عددية أن لفين مغزليين يملكان حالات نقية أكثر مما يمكن أن يوفره لفان مغزليان محضَّر كلٌّ منهما على حدة. وهذه الحالات الزائدة متشابكة، وأبسط مثال عليها يمتلك بالفعل ارتباطات زوجية محددة، في حين أن كل إلكترون على انفراد غير مستقطب تمامًا.}

\section{استكمال مسألة اللف المغزلي الواحد}

يمكن أن تنتظر مبرهنة Bell. يبدأ التشابك نفسه بإلكترونين، لكن البنية المتبقية لإلكترون واحد ستدخل في كل حساب. فلنستكمل إذن النظام ذي الحالتين قبل أن نضم نسختين منه.

اختر متجه وحدة مكاني عاديًا
\begin{equation}
  \hat n=(n_1,n_2,n_3),
  \qquad
  n_1^2+n_2^2+n_3^2=1.
  \label{eq:l04-unit-pointer}
\end{equation}
ويمكن أيضًا أن تُسمّى العلامات \(1,2,3\) بـ\(x,y,z\). وما يهم هو أن المحاور تكوّن جهاز إحداثيات يمينيًا واحدًا متسقًا.

\begin{classroomqa}
\lecturetimestamp{00:02:33}
\audiencequestion{هل الثلاثي الإحداثي المرسوم على اللوح يميني؟}
\lecturerresponse{الرسم التخطيطي الأول أعسري. وقلب المحور الثاني، أو إعادة تسمية محورين على نحو متسق، يصلحه. ولا يستخدم الحساب أدناه إلا اصطلاحًا يمينيًا ثابتًا؛ أما الاتجاه العارض للرسم فلا يحمل أي مضمون فيزيائي.}
\end{classroomqa}

لللف المغزلي ثلاثة مركبات قابلة للرصد، تمثلها مصفوفات باولي
\begin{equation}
  \begin{aligned}
    \sigma_1&=
    \begin{pmatrix}0&1\\1&0\end{pmatrix},
    &
    \sigma_2&=
    \begin{pmatrix}0&-i\\i&0\end{pmatrix},\\[0.5em]
    \sigma_3&=
    \begin{pmatrix}1&0\\0&-1\end{pmatrix}.
  \end{aligned}
\end{equation}
والمركبة على امتداد \(\hat n\) هي الجداء النقطي المكاني المعتاد للمؤشر مع هذا المتجه ذي القيم المصفوفية:
\begin{equation}
  \sigma_{\hat n}
  \equiv
  \boldsymbol{\sigma}\cdot\hat n
  =
  n_1\sigma_1+n_2\sigma_2+n_3\sigma_3.
  \label{eq:l04-sigma-dot-n}
\end{equation}
هذا الجداء النقطي يقع في الفضاء الفيزيائي ثلاثي الأبعاد. وليس هو الجداء الداخلي بين متجهي حالة.

\begin{figure}[htbp]
\centering
\includegraphics[width=0.82\linewidth]{lecture_04_figure_02.png}
\caption{المركبة المكانية على امتداد \(\hat n\) تصبح تركيبًا خطيًا من ثلاثة مرصودات مصفوفية.}
\lectureframe{00:05:09}
\end{figure}

يعطي التعويض بمصفوفات باولي
\begin{equation}
  \sigma_{\hat n}
  =
  \begin{pmatrix}
    n_3 & n_1-in_2\\
    n_1+in_2 & -n_3
  \end{pmatrix}.
  \label{eq:l04-sigma-n-matrix}
\end{equation}
ومن الملائم أن نختصر
\begin{equation}
  n_-=n_1-in_2,
  \qquad
  n_+=n_1+in_2=n_-^*,
\end{equation}
بحيث
\begin{equation}
  \sigma_{\hat n}
  =
  \begin{pmatrix}
    n_3&n_-\\
    n_+&-n_3
  \end{pmatrix}.
\end{equation}
فالعناصر القطرية الحقيقية والعناصر خارج القطرية المترافقة تجعل المؤثر هيرميتيًا.

\begin{figure}[htbp]
\centering
\includegraphics[width=0.84\linewidth]{lecture_04_figure_03.png}
\caption{مرصود اللف المغزلي على محور اعتباطي في صورة صريحة من رتبة اثنين في اثنين.}
\lectureframe{00:07:34}
\end{figure}

عمليًا، يعني قياس \(\sigma_{\hat n}\) توجيه مجال مغناطيسي قوي على امتداد \(\hat n\) ورصد الاستجابة المعتمدة على اللف المغزلي، ومُثِّل ذلك في المحاضرة بما إذا كان فوتون ينبعث. فالمصفوفة والسؤال المختبري وصفان للقياس نفسه.

ستتكفل متطابقتان صغيرتان ببقية الحساب. أولًا،
\begin{equation}
  \sigma_{\hat n}^{\,2}=I.
  \label{eq:l04-sigma-n-square}
\end{equation}
وينتج هذا إما من ضرب المعادلة~\eqref{eq:l04-sigma-n-matrix} في نفسها وإما من استخدام
\begin{equation}
  \sigma_i^2=I,
  \qquad
  \sigma_i\sigma_j=-\sigma_j\sigma_i
  \quad(i\ne j).
\end{equation}
ومن ثم فإن كل قيمة ذاتية \(\lambda\) لـ\(\sigma_{\hat n}\) تحقق \(\lambda^2=1\)، لذا فالنتيجتان الممكنتان الوحيدتان هما
\begin{equation}
  \lambda=\pm1.
\end{equation}
وثانيًا،
\begin{equation}
  n_+n_-
  =n_1^2+n_2^2
  =1-n_3^2
  =(1-n_3)(1+n_3).
  \label{eq:l04-nplus-nminus}
\end{equation}
ليست هذه مبرهنة عميقة، لكنها المتطابقة التي تُبقي جبر المتجهات الذاتية تحت السيطرة.

\section{من التحضير إلى الاحتمال}

حضِّر الإلكترون بحيث يكون لفه المغزلي \(+1\) على امتداد \(\hat n\)، ثم قِس على امتداد اتجاه مختلف \(\hat m\). والسؤال المركزي في مسألة اللف المغزلي الواحد هو
\begin{equation}
  P(m,+\mid n,+)
  =
  P(\sigma_{\hat m}=+1\mid\sigma_{\hat n}=+1).
  \label{eq:l04-central-probability}
\end{equation}
ويقترح التناظر أن الجواب ينبغي أن يعتمد فقط على الزاوية بين \(\hat n\) و\(\hat m\)، وأن تبديل الاتجاهين لا ينبغي أن يغيره. وهذه التوقعات اختبارات مفيدة، لا بدائل عن مسلمة الاحتمال.

\begin{classroomqa}
\lecturetimestamp{00:14:03}
\audiencequestion{هل تخبرنا ميكانيكا الكم نفسها أن مصفوفات باولي الثلاث هي مؤثرات اللف المغزلي، أم إن هذا يُتعلَّم من التجربة؟}
\lecturerresponse{هذا التعيين جزء من البنية التجريبية المتراكمة لميكانيكا الكم، وليس نتيجة تجربة معزولة واحدة. فالتجارب المتكررة على اللف المغزلي تكشف انتظامات احتمالية، وقد بُنيت الصياغة المصفوفية، عبر أعمال Einstein وBorn وHeisenberg وDirac وغيرهم، لترميز تلك الانتظامات. والمسلمات تخمينات ناجحة على نحو استثنائي تقيدها المعطيات.}
\end{classroomqa}

لتكن
\begin{equation}
  \sigma_{\hat n}\ket{n,+}=\ket{n,+},
  \qquad
  \sigma_{\hat m}\ket{m,+}=\ket{m,+}.
  \label{eq:l04-nm-eigenstates}
\end{equation}
والرمزان \(n,+\) و\(m,+\) علامتين داخل رموز الكيت. وهما يحددان الحالتين الذاتيتين اللتين قيمتهما الذاتية هي \(+1\).

\begin{figure}[htbp]
\centering
\includegraphics[width=0.76\linewidth]{lecture_04_figure_04.png}
\caption{الحالتان الذاتيتان الموافقـتان للقيمة \(+1\) للَّف المغزلي على امتداد \(\hat n\) و\(\hat m\).}
\lectureframe{00:17:36}
\end{figure}

أول مقدار ينبغي حسابه هو سعة الاحتمال
\begin{equation}
  \mathcal A(m,+\mid n,+)=\braket{m,+}{n,+}.
\end{equation}
والاحتمال هو مربع مقدارها المطلق:
\begin{equation}
  P(m,+\mid n,+)
  =
  \left|\braket{m,+}{n,+}\right|^2.
  \label{eq:l04-born-nm}
\end{equation}
السعة أولًا، والاحتمال ثانيًا: هذا الترتيب هو الخطوة الكمية الجوهرية.

\begin{figure}[htbp]
\centering
\includegraphics[width=0.78\linewidth]{lecture_04_figure_05.png}
\caption{التراكب بين الحالة \(n\) المحضَّرة والحالة \(m\) المختبرة هو سعة الاحتمال.}
\lectureframe{00:19:59}
\end{figure}

\subsection{إيجاد الحالة الذاتية}

اكتب متجهًا ذاتيًا مرشحًا بقيمة \(+1\) لـ\(\sigma_{\hat n}\) على الصورة
\begin{equation}
  \begin{pmatrix}\alpha\\\beta\end{pmatrix}.
\end{equation}
يمكن تأجيل التطبيع، ولا أثر فيزيائيًا لطور كلي. لذا، بعيدًا عن الرقعة الإحداثية للقطب الجنوبي، يجوز لنا اختيار \(\alpha=1\) وكتابة \(\beta=z\). وتعطي المركبة العلوية من
\begin{equation}
  \begin{pmatrix}
    n_3&n_-\\
    n_+&-n_3
  \end{pmatrix}
  \begin{pmatrix}1\\z\end{pmatrix}
  =
  \begin{pmatrix}1\\z\end{pmatrix}
\end{equation}
\begin{equation}
  n_3+n_-z=1,
  \qquad
  z=\frac{1-n_3}{n_-}.
  \label{eq:l04-z-solution}
\end{equation}
أما المركبة السفلى فتعطي الشرط نفسه بعد استخدام المعادلة~\eqref{eq:l04-nplus-nminus}.

للمتجه غير المطبع مربع معيار يساوي
\begin{align}
  \begin{pmatrix}
    1&\dfrac{1-n_3}{n_+}
  \end{pmatrix}
  \begin{pmatrix}
    1\\[3pt]\dfrac{1-n_3}{n_-}
  \end{pmatrix}
  &=
  1+\frac{(1-n_3)^2}{n_+n_-}
  \notag\\
  &=
  1+\frac{1-n_3}{1+n_3}
  =
  \frac{2}{1+n_3}.
  \label{eq:l04-unnormalized-norm}
\end{align}
ومن ثم فإن اختيارًا مطبَّعًا هو
\begin{equation}
  \ket{n,+}
  =
  \sqrt{\frac{1+n_3}{2}}
  \begin{pmatrix}
    1\\[5pt]\dfrac{1-n_3}{n_-}
  \end{pmatrix}.
  \label{eq:l04-normalized-n-state}
\end{equation}
وعند \(n_3=-1\) تكون هذه الصيغة الإحداثية مفردة مع أن الحالة ليست كذلك؛ فالحالة الذاتية ببساطة هي \(\ket d=(0,1)^T\). وهذا التفرد يعود إلى اصطلاح الطور، لا إلى اللف المغزلي.

واستبدال \(n\) بـ\(m\) ثم أخذ المرافق يعطي
\begin{equation}
  \bra{m,+}
  =
  \sqrt{\frac{1+m_3}{2}}
  \begin{pmatrix}
    1&
    \dfrac{1-m_3}{m_+}
  \end{pmatrix}.
\end{equation}
يبدو التراكب غير مستساغ، ويبدو ضربه في مرافقه أسوأ. لكن الاستخدام المتكرر للمعادلة~\eqref{eq:l04-nplus-nminus} يختزل النتيجة إلى
\begin{equation}
  \begin{aligned}
    P(m,+\mid n,+)
    &=\frac{1+n_1m_1+n_2m_2+n_3m_3}{2}\\
    &=\frac{1+\hat n\cdot\hat m}{2}.
  \end{aligned}
  \label{eq:l04-spin-transition-probability}
\end{equation}
وإذا كانت \(\theta_{nm}\) هي الزاوية بين الاتجاهين المكانيين، فإن
\begin{equation}
  P(m,+\mid n,+)
  =
  \frac{1+\cos\theta_{nm}}{2}
  =
  \cos^2\frac{\theta_{nm}}{2}.
  \label{eq:l04-half-angle-probability}
\end{equation}

\begin{classroomqa}
\lecturetimestamp{00:32:47}
\audiencequestion{هل يجب أن ينعدم التراكب لأن الحالتين متعامدتان؟}
\lecturerresponse{لا. تشير التسميتان \(m,+\) و\(n,+\) إلى محوري قياس مختلفين، لا إلى قيمتين ذاتيتين مختلفتين لمرصود ثابت واحد. فإذا كان \(m=n\)، تطابقت الحالتان وكان تراكبهما واحدًا. ولا تصبحان متعامدتين إلا عندما يكون الاتجاهان الفيزيائيان متعاكسين، بحيث تكون إحداهما حالة \(+1\) والأخرى حالة \(-1\) للمركبة نفسها من اللف المغزلي.}
\end{classroomqa}

\section{ما الذي تعنيه صيغة الاحتمال}

تجتاز المعادلة~\eqref{eq:l04-spin-transition-probability} ثلاثة اختبارات مباشرة:
\begin{align}
  \hat m=\hat n
  &\quad\Longrightarrow\quad P=1,
  \notag\\
  \hat m=-\hat n
  &\quad\Longrightarrow\quad P=0,
  \notag\\
  \hat m\cdot\hat n=0
  &\quad\Longrightarrow\quad P=\frac12.
\end{align}
فالجواب يعتمد فقط على الزاوية النسبية، وهو متناظر تحت \(m\leftrightarrow n\). وتدوير الجهاز كله لا يغير أي تنبؤ.

هذه هي التسوية اللافتة في حالة اللف المغزلي ذي القيمة نصف. فعلى امتداد كل محور تبقى النتيجة المرصودة \(+1\) أو \(-1\)، ولا تكون أبدًا قيمة كسرية للَّف المغزلي. وإنما يظهر الاعتماد الدوراني السلس في احتمالي هذين الجوابين المنفصلين.

\begin{classroomqa}
\lecturetimestamp{00:39:20}
\audiencequestion{كيف نعرف أن المحاولات المتكررة تتعلق بالإلكترون نفسه، وكيف يُقاس الاحتمال؟}
\lecturerresponse{يمكن حبس إلكترون واحد في تجويف ومراقبة الشحنة المحفوظة عبر فيضه الكهربائي، ثم إعادة تحضير ذلك الإلكترون على امتداد \(\hat n\) وتكرار القياس على امتداد \(\hat m\). وبصورة مكافئة يمكن تحضير حشد كبير من الإلكترونات تحضيرًا متماثلًا وقياس كل واحد منها مرة واحدة. والنسبة طويلة الأمد لنتائج الموجب والسالب هي نفسها في الإجرائين، ضمن الخطأ الإحصائي.}
\end{classroomqa}

\begin{classroomqa}
\lecturetimestamp{00:41:29}
\audiencequestion{هل يجب أن تُقاس جميع عناصر الحشد في تجربة جماعية واحدة؟}
\lecturerresponse{لا. يُحضَّر كل عنصر بالبروتوكول نفسه ويُقاس على حدة. ومليون محاولة من هذا النوع تشكل تجربة إحصائية واحدة فقط بالمعنى الذي تكون فيه تواترات نتائجها تقديرًا للاحتمال نفسه.}
\end{classroomqa}

\begin{classroomqa}
\lecturetimestamp{00:42:23}
\audiencequestion{إذا أُعيد استخدام إلكترون واحد، فكم فوتونًا يمكن أن يصدر؟}
\lecturerresponse{لا توجد حصة عمرية ثابتة. تتضمن كل دورة شغلًا خارجيًا عند تشغيل المجال المغناطيسي. وإذا كان اللف المغزلي في الوضع الأعلى طاقة، أمكن لتلك الطاقة المزوَّدة أن تغادر على هيئة فوتون أثناء الاسترخاء. ولذلك يستطيع الجهاز من حيث المبدأ تكرار دورة التحضير والانبعاث إلى ما لا نهاية.}
\end{classroomqa}

يمكن الحصول على الصيغة بسرعة كبيرة بعد معرفة الثبات الدوراني: اختر \(\hat n\) بوصفه المحور الثالث، بحيث \(\ket{n,+}=(1,0)^T\)، واقرأ الاحتمال من المركبة العلوية لـ\(\ket{m,+}\). لكن استعمال هذا الاختصار في البداية كان سيفترض التناظر الذي برهنه الحساب الصريح للتو.

\begin{classroomqa}
\lecturetimestamp{00:43:33}
\audiencequestion{هل يمكن لاختيار أفضل للأساس أن يزيل تلك المتجهات الذاتية القبيحة للمحور الاعتباطي؟}
\lecturerresponse{نعم، ما إن تصبح الدورانات والتحويلات الوحدوية متاحة، يمكن محاذاة أحد المحاور الإحداثية مع أحد الاتجاهين وجعل الحساب قصيرًا. في هذه المرحلة يخدم الحساب الصريح غرضًا: فهو يتحقق من أن مصفوفات باولي المختارة تنتج حقًا جوابًا ثابتًا دورانيًا بدلًا من افتراض تلك الحقيقة.}
\end{classroomqa}

\begin{classroomqa}
\lecturetimestamp{00:48:30}
\audiencequestion{ما السبب الفيزيائي الذي يجعل الجواب يعتمد على الزاوية فقط؟}
\lecturerresponse{ينتمي الجواب البنيوي إلى النظرية الكمومية للتناظر: فالدورانات المكانية تُمثَّل بتحويلات وحدوية تؤثر في حالات اللفّ المغزلي والمرصودات. وستأتي هذه اللغة لاحقًا. أما الآن، فاعتماد النتيجة على الزاوية هو حصيلة مستخلصة من حساب مصفوفات باولي.}
\end{classroomqa}

\section{القياس مرشِّح}

القياس يفعل أكثر من مجرد الإبلاغ عن عدد. فإذا قيس \(\sigma_{\hat m}\) وكانت النتيجة \(+1\)، تُرك الإلكترون في \(\ket{m,+}\). ويكون الجهاز قد عمل ككاشف ومُحضِّر للحالة معًا، على نحو يشبه ما يفعله المستقطِب حين يختبر فوتونًا ويترك الفوتون النافذ مستقطَبًا على امتداد محوره.

هذه القاعدة الخاصة بتغيّر الحالة تجعل مسألة القياس المتزامن ملموسة. ولا يمكن لمرصودين أن تكون لهما معًا قيم محددة قابلة لإعادة الإنتاج إلا عندما توجد حالات هي متجهات ذاتية لكليهما.

\begin{classroomqa}
\lecturetimestamp{00:51:29}
\audiencequestion{هل يمكن لقياس واحد أن يحدد مركبتين مختلفتين للَّف المغزلي، مثل \(\sigma_1\) و\(\sigma_2\)، في آنٍ واحد؟}
\lecturerresponse{لا. فمتجهاتهما الذاتية مختلفة. فقياس إحدى المركبتين يترك الإلكترون في حالة ذاتية لتلك المركبة، لكن تلك الحالة ليست حالة ذاتية للأخرى. ولذلك يدمّر القياس الأول التعيّن اللازم للثانية.}
\end{classroomqa}

تُصاغ العبارة العامة في الفضاءات المنتهية الأبعاد بأوضح صورة لمؤثرات هرميتية.

\begin{theorem}
يمكن قطرنة مؤثرين هرمِيتيين \(A\) و\(B\) في آن واحد إذا وفقط إذا كانا متبادلين:
\begin{equation}
  [A,B]\equiv AB-BA=0.
  \label{eq:l04-commutator}
\end{equation}
وعلى نحو مكافئ، يملكان أساسًا ذاتيًا متعامدًا مُطبَّعًا كاملًا مشتركًا.
\end{theorem}

في الاتجاه الأسهل، افترض أن
\begin{equation}
  A\ket{\alpha_i}=\lambda_i\ket{\alpha_i},
  \qquad
  B\ket{\alpha_i}=\mu_i\ket{\alpha_i}
\end{equation}
بالنسبة إلى أساس كامل. وعندئذٍ
\begin{align}
  AB\ket{\alpha_i}
  &=\lambda_i\mu_i\ket{\alpha_i},
  \notag\\
  BA\ket{\alpha_i}
  &=\mu_i\lambda_i\ket{\alpha_i},
\end{align}
ومن ثم فإن \(AB\) و\(BA\) يتطابقان على كل متجه أساس. إذن \(AB=BA\). وبالعكس، إذا كان \(A\) و\(B\) هرمِيتيين ومتبادلين، فإن \(B\) يحفظ كل فضاء ذاتي لـ\(A\)؛ وقطرنة \(B\) داخل أي فضاء ذاتي منحلّ تنتج أساسًا ذاتيًا مشتركًا.

\begin{figure}[htbp]
\centering
\includegraphics[width=0.76\linewidth]{lecture_04_figure_06.png}
\caption{قد يكون لمرصودين المتجهات الذاتية نفسها مع إسناد قيم ذاتية مختلفة إليها.}
\lectureframe{00:55:45}
\end{figure}

\begin{classroomqa}
\lecturetimestamp{00:53:56}
\audiencequestion{هل يمكننا تطبيق مجالين مغناطيسيين على امتداد \(x\) و\(y\) معًا ونسمي ذلك قياسًا متزامنًا؟}
\lecturerresponse{تتجمع الحقول في اتجاه محصِّل واحد. وعندئذٍ يقيس الجهاز اللفّ المغزلي على امتداد ذلك المحور المحصِّل، ربما عند \(45^\circ\)، بدلًا من قياس \(\sigma_x\) و\(\sigma_y\) كلٌّ على حدة. فالبناء المُحاوَل يُنتج مرصودًا ثالثًا، لا قراءتين متوافقتين.}
\end{classroomqa}

\section{لكل حالة نقية أحادية اللفّ المغزلي محور}

والآن اعكس مسألة المتجه الذاتي. إذا أُعطي سبينور مُطبَّع اعتباطي
\begin{equation}
  \ket\psi=
  \begin{pmatrix}\alpha\\\beta\end{pmatrix},
  \qquad
  |\alpha|^2+|\beta|^2=1,
  \label{eq:l04-general-spinor}
\end{equation}
فهل يوجد اتجاه مكاني \(\hat n\) بحيث يحقق
\begin{equation}
  \sigma_{\hat n}\ket\psi=\ket\psi?
  \label{eq:l04-reverse-eigenproblem}
\end{equation}
نعم. وترى المحاضرة أولًا سبب ذلك عن طريق العدّ. فالمركبتان العقديتان تحتويان أربعة أعداد حقيقية. يزيل التطبيع واحدًا منها، ويزيل الطور الكلي غير ذي الأهمية الفيزيائية واحدًا آخر، فيبقى اثنان. كما أن مؤشرًا وحديًا له أيضًا معلمتان حقيقتان مستقلتان، مثل خط الطول وخط العرض.

\begin{classroomqa}
\lecturetimestamp{01:03:52}
\audiencequestion{كيف يمكن أن تكون للحالة درجتا حرية فقط مع أن \(\alpha\) و\(\beta\) كلاهما عقدي؟}
\lecturerresponse{تبدأ بأربع مركبات حقيقية. يزيل شرط التطبيع درجة حرية حقيقية واحدة، وضرب الحالة كلها في \(e^{i\theta}\) لا يغيّر أي احتمال، لذا يزيل الطور المشترك درجة أخرى. وتبقى معلمتان ذواتا معنى فيزيائي، وتتطابقان مع الزاويتين اللتين تحددان مؤشرًا وحديًا.}
\end{classroomqa}

\begin{classroomqa}
\lecturetimestamp{01:04:20}
\audiencequestion{هل نحتاج أيضًا إلى إيجاد القيمة الذاتية؟}
\lecturerresponse{هنا تُثبَّت القيمة الذاتية المستهدفة عند \(+1\). ونحن لا نبحث إلا عن الاتجاه الذي تكون حالته لأعلى هي السبينور المعطى. أما المؤشر المعاكس فسيصف الشعاع نفسه بوصفه حالة لأسفل ذات قيمة ذاتية \(-1\).}
\end{classroomqa}

يشرح العدّ لماذا ينبغي أن توجد هذه المطابقة. ويؤكدها إنشاء مباشر. عرّف
\begin{align}
  n_1&=\bra\psi\sigma_1\ket\psi
      =2\operatorname{Re}(\alpha^*\beta),
  \notag\\
  n_2&=\bra\psi\sigma_2\ket\psi
      =2\operatorname{Im}(\alpha^*\beta),
  \notag\\
  n_3&=\bra\psi\sigma_3\ket\psi
      =|\alpha|^2-|\beta|^2.
  \label{eq:l04-bloch-components}
\end{align}
يعطي التطبيع \(n_1^2+n_2^2+n_3^2=1\)، ويعطي التعويض المباشر المعادلة ~\eqref{eq:l04-reverse-eigenproblem}. ومن ثم فكل حالة نقية أحادية اللفّ المغزلي هي حالة لأعلى على امتداد محور مكاني ما.

وثمة صيغة إجرائية لهذه العبارة. حضِّر نسخًا كثيرة من الحالة النقية المجهولة نفسها. قِس دفعاتٍ على امتداد اتجاهات تجريبية، واضبط المغناطيس، وابحث عن الاتجاه الذي تكون فيه كل النتائج لأعلى. إن البيانات المنتهية لا تحدد الاتجاه إلا في حدود عدم يقين، لكن للحالة المثالية محورًا دقيقًا واحدًا.

الأسئلة جزء من هذه الطريقة، لا مقاطعات لها. وتستعيد استطرادة صفية قصيرة أطفالًا في العاشرة من العمر يسألون بحرية عن نظرية الأوتار، والأبعاد الإضافية، وما إذا كانت النيوترينوهات المحيطة تسبب الصداع. وقد تكون تفاصيل السؤال خاطئة؛ لكن الاستعداد لكشف موضع الالتباس يظل ذا قيمة.

\section{لفّان مغزليان ومؤثرات موضعية}

قبل جمع لفّين مغزليين، سجِّل كيف تؤثر مصفوفات باولي في أساس \(\sigma_3\):
\begin{align}
  \sigma_1\ket u&=\ket d,
  &
  \sigma_1\ket d&=\ket u,
  \notag\\
  \sigma_2\ket u&=i\ket d,
  &
  \sigma_2\ket d&=-i\ket u,
  \notag\\
  \sigma_3\ket u&=\ket u,
  &
  \sigma_3\ket d&=-\ket d.
  \label{eq:l04-pauli-action-table}
\end{align}

\begin{figure}[htbp]
\centering
\includegraphics[width=0.80\linewidth]{lecture_04_figure_07.png}
\caption{يعيد جدول تأثير باولي صياغة المصفوفات بوصفها قواعد على أساس لأعلى/لأسفل.}
\lectureframe{01:09:35}
\end{figure}

تخيّل الآن إلكترونين موضوعين في موضعين متميزين بحيث يكون المتغير الوحيد هو لفّهما المغزلي. حالات الأساس الأربع هي
\begin{equation}
  \ket{uu},
  \qquad
  \ket{du},
  \qquad
  \ket{ud},
  \qquad
  \ket{dd}.
  \label{eq:l04-two-spin-basis}
\end{equation}
الوسم الأول يخص الإلكترون \(1\)، والثاني يخص الإلكترون \(2\). وهذه الحالات متعامدة مُطبَّعة فيما بينها لأنها تمثل نواتج مشتركة متميزة لقياسي المركبة الثالثة للجسيمين.

\begin{figure}[htbp]
\centering
\includegraphics[width=0.82\linewidth]{lecture_04_figure_08.png}
\caption{تولد وسوم أساس الجداء الأربعة فضاء حالات اللفّين المغزليين.}
\lectureframe{01:11:40}
\end{figure}

سمِّ المؤثرات المؤثرة في الإلكترون \(1\) بـ\(\sigma_i\)، والمؤثرات المؤثرة في الإلكترون \(2\) بـ\(\tau_i\). وبصورة أكثر صراحةً،
\begin{equation}
  \sigma_i\equiv\sigma_i\otimes I,
  \qquad
  \tau_i\equiv I\otimes\sigma_i.
\end{equation}
تؤثر المجموعة الأولى في الوسم الأول فقط، بينما تؤثر الثانية في الوسم الثاني فقط:
\begin{align}
  \sigma_1\ket{ud}&=\ket{dd},
  &
  \sigma_1\ket{dd}&=\ket{ud},
  \notag\\
  \sigma_3\ket{dd}&=-\ket{dd},
  &
  \tau_2\ket{ud}&=-i\ket{uu}.
  \label{eq:l04-local-operator-examples}
\end{align}
ومتى عُرف فعل مؤثر ما على حالات الأساس الأربع، حددت الخطية فعله على كل حالة للَّفّين المغزليين.

\begin{classroomqa}
\lecturetimestamp{01:17:47}
\audiencequestion{لماذا كُتب العامل \(-i\) خارج الكيت؟ هل يؤثر في الإلكترونين كليهما؟}
\lecturerresponse{الرموز داخل الكيت هي وسوم لتهيئة من تهيئات الأساس. والمقدار العددي العقدي يضرب متجه الحالة كله، لذا تضعه الكتابة القياسية خارج الكيت. ومصدره هنا هو فعل \(\tau_2\) على الوسم الثاني؛ وكتابته في الخارج لا تعني أن عملية فيزيائية منفصلة قد طُبِّقت على الإلكترونين كليهما.}
\end{classroomqa}

\begin{classroomqa}
\lecturetimestamp{01:19:44}
\audiencequestion{هل يمثّل الضرب في \(i\) دورانًا للإلكترون في الفضاء العادي؟}
\lecturerresponse{ليس بمفرده. إنه ضرب بعدد عقدي في فضاء الحالة. وستُمثَّل الدورانات المكانية لاحقًا بمؤثرات وحدوية مخصوصة؛ أما العامل المنفرد \(i\) في جدول الفعل هذا فليس دورانًا حرفيًا في الفضاء الحقيقي.}
\end{classroomqa}

\section{ستّ معلمات لا أربع}

الحالة النقية العامة للزوج هي
\begin{equation}
  \ket\Psi
  =
  c_{uu}\ket{uu}
  +c_{du}\ket{du}
  +c_{ud}\ket{ud}
  +c_{dd}\ket{dd}.
  \label{eq:l04-general-two-spin-state}
\end{equation}
تعطي المعاملات العقدية الأربعة ثمانية أعداد حقيقية. يزيل التطبيع واحدًا ويزيل الطور الكلي واحدًا:
\begin{equation}
  8-1-1=6.
  \label{eq:l04-two-spin-parameter-count}
\end{equation}
وتحتاج حالة نقية مُحضَّرة على نحو مستقل للإلكترون \(1\) إلى معلمتين حقيقتين، ويحتاج الإلكترون \(2\) إلى معلمتين أخريين، فيكون المجموع أربعة فقط. أما المعلمتان الناقصتان فليستا خطأً في العدّ. إنهما تدلان على حالات نقية للزوج لا يمكن إسناد حالات نقية منفصلة فيها إلى الجزأين.

لنحدِّد أولًا الحالات التي تأتي فعلًا من تحضيرات منفصلة. إذا
\begin{equation}
  \ket{\psi_1}=\alpha_1\ket u+\beta_1\ket d,
  \qquad
  \ket{\psi_2}=\alpha_2\ket u+\beta_2\ket d,
\end{equation}
فعندئذٍ
\begin{align}
  \ket{\psi_1}\otimes\ket{\psi_2}
  &=
  \alpha_1\alpha_2\ket{uu}
  +\alpha_1\beta_2\ket{ud}
  \notag\\
  &\quad
  +\beta_1\alpha_2\ket{du}
  +\beta_1\beta_2\ket{dd}.
  \label{eq:l04-product-expansion}
\end{align}
فهذه \emph{حالة جداء}. ولكل إلكترون محور استقطابه الخاص، كما لو أن التحضيرين المغناطيسيين أُجريا على نحو مستقل.

\begin{figure}[htbp]
\centering
\includegraphics[width=0.84\linewidth]{lecture_04_figure_09.png}
\caption{إن ضرب حالتين أحاديتَي اللفّ المغزلي ينتج حالة جداء مقيَّدة ذات أربعة حدود.}
\lectureframe{01:25:15}
\end{figure}

إذا رُتِّبت المعاملات الأربعة على الصورة
\begin{equation}
  C=
  \begin{pmatrix}
    c_{uu}&c_{ud}\\
    c_{du}&c_{dd}
  \end{pmatrix},
\end{equation}
فإن الحالة تكون حالة جداء إذا وفقط إذا كانت رتبة \(C\) تساوي واحدًا، وبصورة مكافئة
\begin{equation}
  c_{uu}c_{dd}-c_{ud}c_{du}=0.
  \label{eq:l04-product-condition}
\end{equation}
ومعظم اختيارات السعات الأربع لا تحقق هذا الشرط. فهي متشابكة.

\section{أول حالة متشابكة}

أول مثال في المحاضرة هو
\begin{equation}
  \ket{\psi_+}
  =
  \frac{1}{\sqrt2}
  \left(\ket{ud}+\ket{du}\right).
  \label{eq:l04-psi-plus}
\end{equation}
ولا يمكن تحليلها إلى حالة للإلكترون \(1\) وأخرى للإلكترون \(2\). فإذا وجد قياس على المحور الثالث الإلكترون \(2\) إلى أسفل، كان الإلكترون \(1\) إلى أعلى؛ وإذا وجد القياس الإلكترون \(2\) إلى أعلى، كان الإلكترون \(1\) إلى أسفل. وللزوج معلومات محددة لا يمتلكها أي من الجسيمين على انفراد.

\begin{figure}[htbp]
\centering
\includegraphics[width=0.82\linewidth]{lecture_04_figure_10.png}
\caption{أول تراكب متشابك يحتوي على واحدٍ لأعلى وواحدٍ لأسفل، لكنه لا يعيّن أيًّا من النتيجتين لجسيم بعينه سلفًا.}
\lectureframe{01:29:30}
\end{figure}

لكي نرى كيف يبدو أيٌّ من الإلكترونين على حدة، نحسب قيم التوقع الموضعية الثلاث. بالنسبة إلى المركبة الأولى،
\begin{equation}
  \sigma_1\ket{\psi_+}
  =
  \frac{1}{\sqrt2}
  \left(\ket{dd}+\ket{uu}\right),
\end{equation}
وهذا متعامد مع \(\ket{\psi_+}\). لذلك
\begin{equation}
  \bra{\psi_+}\sigma_1\ket{\psi_+}=0.
\end{equation}
وبالمثل،
\begin{equation}
  \sigma_2\ket{\psi_+}
  =
  \frac{i}{\sqrt2}
  \left(\ket{dd}-\ket{uu}\right),
\end{equation}
ومن ثم
\begin{equation}
  \bra{\psi_+}\sigma_2\ket{\psi_+}=0.
\end{equation}
أما بالنسبة إلى المركبة الثالثة،
\begin{equation}
  \sigma_3\ket{\psi_+}
  =
  \frac{1}{\sqrt2}
  \left(\ket{ud}-\ket{du}\right),
\end{equation}
ويلغي الإسهامان القطريان أحدهما الآخر:
\begin{equation}
  \bra{\psi_+}\sigma_3\ket{\psi_+}
  =\frac12(1-1)=0.
\end{equation}
وعليه
\begin{equation}
  \begin{aligned}
    \bra{\psi_+}\sigma_i\ket{\psi_+}&=0,\\
    \bra{\psi_+}\tau_i\ket{\psi_+}&=0,
    \qquad i=1,2,3.
  \end{aligned}
  \label{eq:l04-zero-local-spin}
\end{equation}

\begin{figure}[htbp]
\centering
\includegraphics[width=0.84\linewidth]{lecture_04_figure_11.png}
\caption{تثبت الحالة المتحوِّلة والإلغاء \(1-1=0\) انعدام قيم توقّع اللف المغزلي الموضعية.}
\lectureframe{01:35:50}
\end{figure}

بالنسبة إلى أي اتجاه،
\begin{equation}
  \bra{\psi_+}\sigma_{\hat n}\ket{\psi_+}=0,
  \qquad
  \bra{\psi_+}\tau_{\hat n}\ket{\psi_+}=0.
  \label{eq:l04-zero-arbitrary-local-spin}
\end{equation}
وبما أن لكل قياس موضعي ناتجين فقط هما \(\pm1\)، فإن انعدام قيمة التوقع يعني
\begin{equation}
  P(\sigma_{\hat n}=+1)
  =
  P(\sigma_{\hat n}=-1)
  =
  \frac12,
\end{equation}
وكذلك بالنسبة إلى الإلكترون الثاني. فلا يوجد اتجاه يستقطب أيًّا من الإلكترونين بمفرده. وهذا مستحيل لحالة نقية لسبين واحد، لكنه ممكن تمامًا لأحد جزأي حالة نقية متشابكة.

\begin{classroomqa}
\lecturetimestamp{01:38:23}
\audiencequestion{هل النتيجة الموضعية نصفًا بنصف، وكون قياس أحد الإلكترونين يزوّدنا بمعلومة عن الآخر، هما الادعاء نفسه؟}
\lecturerresponse{لا. تجاهل أولًا الإلكترون \(2\): فكل اتجاه يُقاس على الإلكترون \(1\) وحده يعطي توزيعًا نصفًا بنصف. وهذا يثبت غياب أي استقطاب موضعي. أما الارتباط فهو قضية مشتركة منفصلة عن كيفية ارتباط النتيجتين المسجَّلتين إحداهما بالأخرى. ويوفّر التشابك السمتين كلتيهما، لكن لا ينبغي استعمال إحداهما بديلًا عن إثبات الأخرى.}
\end{classroomqa}

\section{الارتباط قبل بيل}

الارتباط وحده ليس كميًّا على نحو فريد. اخلط قطعة من فئة السنت وأخرى من فئة العشرة سنتات، وافصل بينهما من دون نظر، ثم افحص إحداهما. قبل الفحص يكون لكل واحدة منهما احتمال متساوٍ لأن تكون قطعة السنت، لكن العثور على قطعة السنت هنا يخبرنا بأن القطعة البعيدة هي قطعة العشرة سنتات. فالزوج يحمل معلومة محددة حتى حين لا يحملها أي جزء منه بمفرده.

\begin{classroomqa}
\lecturetimestamp{01:41:45}
\audiencequestion{هل يعمل استدلال السنت والعشرة سنتات فقط إذا كنا نعلم أن الزوج الأصلي كان يحتوي بالضبط على واحدة من كل نوع؟}
\lecturerresponse{بالضبط. فالاستدلال الكلاسيكي يعتمد على إعداد موثوق. وكذلك لا ينطبق تنبؤ الارتباط الكمي إلا عندما تكون منظومة الإلكترونين قد أُعِدَّت في الحالة المشتركة المحددة.}
\end{classroomqa}

\begin{classroomqa}
\lecturetimestamp{01:42:09}
\audiencequestion{بالنسبة إلى \(\ket{\psi_+}\)، هل يساوي احتمال أن تكون نتيجتا قياس المركبة الثالثة كلتاهما إلى أعلى الصفر؟ وهل يثبت ذلك أن اللفّين المغزليين متعاكسان على كل محور؟}
\lecturerresponse{احتمال \(uu\) في أساس المحور الثالث يساوي صفرًا، ونتيجتا المحور الثالث متعاكستان. لكن الاستنتاج الثاني أقوى مما تسمح به هذه الحالة ذات الإشارة الموجبة. فاللف المغزلي الموضعي لهذه الحالة يساوي صفرًا في كل اتجاه، بينما يعتمد ارتباطها المشترك على الاتجاه. أما تعاكس النتائج على كل محور فيخص حالة السينغلت ذات الإشارة السالبة التي تُعرض في نهاية المحاضرة.}
\end{classroomqa}

\begin{classroomqa}
\lecturetimestamp{01:43:51}
\audiencequestion{هل يمكن تخيل الزوج الكلاسيكي على أنه قطعتان نقديتان ملتحمتان بحيث يكون أحد الوجهين على الدوام معاكسًا للآخر؟}
\lecturerresponse{هذه آلية كلاسيكية ممكنة لارتباط مضاد ثابت، بشرط ألا يفسد الهزُّ ذلك الالتحام. وهو يظل ترتيبًا كلاسيكيًا خفيًا ثُبِّت عند الإعداد. وهذا التشبيه مؤقت عمدًا؛ فمبرهنة بيل ستختبر موضع إخفاق كل تفسير كلاسيكي من هذا النوع.}
\end{classroomqa}

\begin{classroomqa}
\lecturetimestamp{01:44:32}
\audiencequestion{أين يدخل قانون الأعداد الكبيرة إذا كانت إحدى النتيجتين تحدد بالفعل النتيجة الأخرى؟}
\lecturerresponse{اليقين الشرطي حكم يتعلق بتجربة واحدة متى أُعطيت الحالة المشتركة وإحدى النتيجتين. ومع ذلك تظل التجارب المتكررة ضرورية للتحقق تجريبيًا من أن الارتباط المتنبأ به موثوق، ولتقدير قيم التوقع والاحتمالات التي ليست يقينية. ويربط قانون الأعداد الكبيرة بين تلك التواترات المتكررة والتوزيع النظري.}
\end{classroomqa}

قد يكون المراقبان مفصولين بمسافة هائلة. إن النظر إلى إحدى النتيجتين يغيّر ما يعرفه المراقب المحلي عن النتيجة الأخرى، لكنه لا يرسل رسالة يمكن التحكم فيها إلى المراقب البعيد. وفي هذه المرحلة لا يزال الوضع يشبه قطعتَي السنت والعشرة سنتات بعد فصلهما.

\begin{classroomqa}
\lecturetimestamp{01:46:46}
\audiencequestion{هل صورة السنت والعشرة سنتات بعد فصلهما هي الطريقة الصحيحة للتفكير في الزوج الكمي؟}
\lecturerresponse{إنها المقارنة الكلاسيكية الأولى الصحيحة، وليست التفسير الكمي النهائي. فمبرهنة بيل تتعلق بالفرق القابل للقياس بين زوج يحمل خواص كلاسيكية موجودة سلفًا وزوج كمي متشابك. ويُؤجَّل هذا التمييز إلى أن تُحسب الارتباطات اللازمة.}
\end{classroomqa}

\section{حالة السينغلت وتحضيرها}

في نهاية الحساب على السبورة، استبدل إشارة الجمع بإشارة طرح:
\begin{equation}
  \ket{\psi_-}
  =
  \frac{1}{\sqrt2}
  \left(\ket{ud}-\ket{du}\right).
  \label{eq:l04-singlet}
\end{equation}
يبقي هذا التغيير قيم التوقع الموضعية في المعادلة~\eqref{eq:l04-zero-local-spin} مساوية للصفر، ويبقي نتائج المحور الثالث متعاكسة. لكنه يغيّر الارتباطات المشتركة المستعرضة. وهذه الحالة ذات الإشارة السالبة هي سينغلت اللف المغزلي، أي الحالة غير المتغيرة دورانيًا التي يكون فيها اللفّان المغزليان متعاكسين على كل محور قياس مشترك.

\begin{figure}[htbp]
\centering
\includegraphics[width=0.84\linewidth]{lecture_04_figure_12.png}
\caption{يستبدل التغيير الأخير على السبورة التراكب المتماثل بسينغلت ذي إشارة سالبة مع الإبقاء على انعدام اللف المغزلي الموضعي.}
\lectureframe{01:47:25}
\end{figure}

في الصورة الفيزيائية المبسطة المستخدمة هنا، فإن تقريب العزمين المغناطيسيين المتموضعين لإلكترونين بما يكفي لكي يتآثرا يسمح لاصطفاف الطاقة الأعلى بأن يسترخي عبر انبعاث فوتون. وفي غياب محور خارجي مفضَّل، تُمثَّل حالة اللف المغزلي الأدنى بحالة السينغلت المتناظرة دورانيًا. أما التآثر الدقيق وسبب وجوب كون الطور النسبي سالبًا فيتطلبان حسابًا إضافيًا.

\begin{classroomqa}
\lecturetimestamp{01:49:02}
\audiencequestion{هل يكون التحضير العملي ببساطة هو تقريب إلكترونين من بعضهما والانتظار؟}
\lecturerresponse{في هذا النموذج المثالي، نعم: يُترَك تآثر اللف المغزلي-اللف المغزلي ليعمل، ويُسمَح بالاسترخاء الإشعاعي إلى الحالة الأدنى. ويكفي أن يكون الإلكترونان متقاربين بما يكفي لمدة كافية لكي يكون تآثرهما المغناطيسي فعالًا. أما بروتوكول التحضير الحقيقي فعليه أيضًا أن يضبط الحالة المدارية والتآثرات المنافسة.}
\end{classroomqa}

\begin{classroomqa}
\lecturetimestamp{01:49:16}
\audiencequestion{هل هذا مجرد مبدأ باولي للاستبعاد؟}
\lecturerresponse{لا. فقد أُسنِدت إلى الإلكترونين مواضع مكانية متميزة، لذلك فإن مناقشة لفِّهما المغزلي لا تتطلب بعدُ جعل الدوال الموجية المدارية المتراكبة مضادّة للتناظر. ويصبح الاستبعاد ذا صلة حين تشغل إلكترونات متماثلة حالات أحادية الجسيم متراكبة؛ أما حجة التشابك الحالية فتعمل بالفعل للمواقع المنفصلة.}
\end{classroomqa}

\begin{classroomqa}
\lecturetimestamp{01:50:31}
\audiencequestion{هل ستحتوي حزمة إلكترونات عادية طبيعيًا على أزواج متآثرة من هذا النوع؟}
\lecturerresponse{عادةً لا. فالحزمة مخففة، وتبقى الإلكترونات النموذجية متباعدة أكثر مما يسمح بتآثر لف مغزلي-لف مغزلي ملحوظ خلال زمن الطيران المتاح. ويمكن للتصادمات أو لتآثر حابس مقصود أن ينشئا أزواجًا متشابكة، لكن مجرد الانتماء إلى الحزمة نفسها لا يفعل ذلك.}
\end{classroomqa}

\begin{classroomqa}
\lecturetimestamp{01:51:45}
\audiencequestion{ما المسافة المطلوبة للتآثر المغناطيسي؟}
\lecturerresponse{لا يوجد حد فاصل حاد عام؛ فالكمية الحاسمة هي المقياس الزمني للتآثر. وتعطي المحاضرة تقديرًا مجهريًا تقريبيًا من رتبة \(10^{-11}\,\mathrm{cm}\) لاقتران سريع. وعند تباعد أكبر يكون الانقسام الطاقي أصغر، ويكون للفوتون المنبعث طول موجي أطول، ويستغرق الاسترخاء وقتًا أطول بكثير. ومن حيث المبدأ يمكن لانتظار طويل بما يكفي أن يعوّض ذلك، بينما تتيح الحزمة العادية وقتًا أقل بكثير مما يلزم.}
\end{classroomqa}

\begin{classroomqa}
\lecturetimestamp{01:53:10}
\audiencequestion{لماذا يجب أن تحتوي حالة السينغلت على إشارة سالبة بدلًا من إشارة موجبة؟}
\lecturerresponse{هذا ليس واضحًا من الإحصاءات الموضعية نصفًا بنصف، لأن هذه الإحصاءات متطابقة في الطورين النسبيين كليهما. ويكمن الجواب في كيفية تحوّل الحالة تحت الدورانات، وفي كيفية اختيار التآثر للَّف المغزلي الكلي الذي يساوي صفرًا. أما الاشتقاق فينتمي إلى المحاضرة التالية.}
\end{classroomqa}

لقد تجاوز المقرر الآن العتبة المقصودة. فحالة سبين نقية مفردة لها دائمًا محور محدد تكون عليه إلى أعلى. أما الزوج النقي فلا يلزم أن يسند إلى أيٍّ من عضويه حالة نقية، ومع ذلك يمكنه أن يشفّر علاقات دقيقة بين نتائجهما. وهذه البنية الإضافية هي التشابك؛ وستبيّن مبرهنة بيل لماذا لا تستطيع أي قصة كلاسيكية من نوع السنت والعشرة سنتات أن تعيد إنتاج ذلك كله.
